\kchapter{Extending Kubernetes: CRDs, Controllers, Operators}{crd}

\begin{objectives}
\begin{itemize}
  \item how Kubernetes can learn brand-new resource types;
  \item CustomResourceDefinitions (CRDs) and custom resources;
  \item the controller pattern you can implement yourself;
  \item what Operators are and when to use them;
  \item admission webhooks and the wider extension landscape.
\end{itemize}
\end{objectives}

\section{Kubernetes as a platform for platforms}

Everything you have learned follows one pattern: \emph{an object stores desired
state, and a controller reconciles reality toward it}. Kubernetes lets you add
your own objects and your own controllers using exactly the same machinery. This
is why tools like cert-manager, Argo CD, Prometheus Operator, Crossplane, and
KEDA feel ``native'': they are.

\section{CustomResourceDefinitions}

A \kicon{crd}~\kterm{CRD} teaches the API server a new \kc{kind}, complete
with an OpenAPI schema for validation, \kc{kubectl} support, RBAC, and watch
events.

\begin{yamlcode}{website-crd.yaml}
apiVersion: apiextensions.k8s.io/v1
kind: CustomResourceDefinition
metadata:
  name: websites.demo.example.com      # <plural>.<group>
spec:
  group: demo.example.com
  scope: Namespaced
  names:
    kind: Website
    plural: websites
    singular: website
    shortNames: [ws]
  versions:
    - name: v1alpha1
      served: true
      storage: true
      schema:
        openAPIV3Schema:
          type: object
          properties:
            spec:
              type: object
              required: [gitRepo]
              properties:
                gitRepo:  { type: string }
                replicas: { type: integer, minimum: 1, maximum: 10, default: 2 }
            status:
              type: object
              properties:
                url: { type: string }
      subresources:
        status: {}
      additionalPrinterColumns:
        - name: Replicas
          type: integer
          jsonPath: .spec.replicas
        - name: URL
          type: string
          jsonPath: .status.url
\end{yamlcode}

\begin{yamlcode}{my-site.yaml}
apiVersion: demo.example.com/v1alpha1
kind: Website
metadata:
  name: docs
spec:
  gitRepo: https://github.com/example/docs-site
  replicas: 3
\end{yamlcode}

\begin{bashcode}{terminal}
kubectl apply -f website-crd.yaml
kubectl apply -f my-site.yaml
kubectl get websites          # or: kubectl get ws
\end{bashcode}

On its own, a CRD is just storage: the \kc{Website} is saved, but nothing
happens. To make it \emph{do} something you need a controller.

\section{Writing a controller}

A controller watches your custom resources (and the objects it creates) and
reconciles. For the \kc{Website} example, the reconcile function might:

\begin{enumerate}
  \item read the \kc{Website} object;
  \item ensure a Deployment exists that clones \kc{gitRepo} and serves it, with
        \kc{replicas} copies;
  \item ensure a Service and an HTTPRoute exist;
  \item write the public URL into \kc{status.url};
  \item set \kc{ownerReferences} so deleting the Website garbage-collects everything.
\end{enumerate}

\begin{plaincode}{reconcile.go (sketch, using controller-runtime)}
func (r *WebsiteReconciler) Reconcile(ctx context.Context, req ctrl.Request) (ctrl.Result, error) {
    var site demov1alpha1.Website
    if err := r.Get(ctx, req.NamespacedName, &site); err != nil {
        return ctrl.Result{}, client.IgnoreNotFound(err)   // deleted: nothing to do
    }

    deploy := &appsv1.Deployment{ObjectMeta: metav1.ObjectMeta{Name: site.Name, Namespace: site.Namespace}}
    _, err := controllerutil.CreateOrUpdate(ctx, r.Client, deploy, func() error {
        deploy.Spec.Replicas = &site.Spec.Replicas
        // ... fill in the Pod template from site.Spec.GitRepo ...
        return controllerutil.SetControllerReference(&site, deploy, r.Scheme)
    })
    if err != nil {
        return ctrl.Result{}, err                          // retried with back-off
    }

    site.Status.URL = "https://" + site.Name + ".example.com"
    return ctrl.Result{}, r.Status().Update(ctx, &site)
}
\end{plaincode}

\begin{deepdive}
Reconcile functions must be \textbf{idempotent} and \textbf{level-triggered}:
they look at the whole current state each time, not at ``what changed''. If a
reconcile fails or the controller restarts, the next run simply converges again.
Frameworks: \textbf{Kubebuilder} and \textbf{Operator SDK} (Go),
\textbf{kopf} (Python), \textbf{Java Operator SDK}, \textbf{kube-rs} (Rust).
\end{deepdive}

\section{Operators}

An \kterm{Operator} is a controller that encodes the operational knowledge of a
human expert for a specific piece of software: how to install it, upgrade it,
back it up, fail over, and scale it.

\begin{center}
\ktablestart
\begin{tabularx}{\linewidth}{P{4.0cm} L}
\kheader \kth{Operator} & \kth{Manages} \\
CloudNativePG & PostgreSQL clusters with replication, failover, backups to object storage. \\
Strimzi & Apache Kafka clusters, topics, and users. \\
Prometheus Operator & Prometheus, Alertmanager, ServiceMonitors, rules. \\
cert-manager & Certificates and issuers. \\
Crossplane & Cloud resources (databases, buckets, networks) as Kubernetes objects. \\
\end{tabularx}
\end{center}

\begin{yamlcode}{cnpg-cluster.yaml}
apiVersion: postgresql.cnpg.io/v1
kind: Cluster
metadata:
  name: taskflow-db
spec:
  instances: 3                 # 1 primary + 2 replicas, automatic failover
  storage:
    size: 20Gi
  backup:
    barmanObjectStore:
      destinationPath: s3://my-backups/taskflow
      s3Credentials:
        accessKeyId:     { name: s3-creds, key: ACCESS_KEY_ID }
        secretAccessKey: { name: s3-creds, key: SECRET_ACCESS_KEY }
\end{yamlcode}

\begin{keyidea}
Three lines of YAML (\kc{instances: 3}) replace pages of database runbooks.
That is the promise of Operators --- and why many teams run stateful systems on
Kubernetes today.
\end{keyidea}

\section{Admission webhooks}

Besides new types, you can intercept API requests:

\begin{itemize}
  \item \textbf{Mutating} admission webhooks modify objects as they are created
        (inject sidecars, add default labels).
  \item \textbf{Validating} admission webhooks accept or reject them (enforce policy).
\end{itemize}

For many policies, the built-in \textbf{ValidatingAdmissionPolicy} and
\textbf{MutatingAdmissionPolicy} (CEL expressions, no server to run) are now
simpler than webhooks.

\section{Other extension points}

\begin{itemize}
  \item \textbf{CNI}, \textbf{CSI}, \textbf{CRI} --- network, storage, runtime plugins.
  \item \textbf{Device plugins} --- GPUs and other hardware (with Dynamic Resource Allocation in newer releases).
  \item \textbf{Scheduler plugins} --- custom scheduling logic.
  \item \textbf{API aggregation} --- whole extra API servers (metrics-server uses this).
  \item \textbf{kubectl plugins} --- any \kc{kubectl-foo} binary on your PATH
        becomes \kc{kubectl foo}; manage them with \textbf{krew}.
\end{itemize}

\begin{recap}
\begin{itemize}
  \item CRDs add new resource types; controllers give them behaviour.
  \item Operators package expert operational knowledge as controllers.
  \item Admission control lets you mutate or validate every request.
\end{itemize}
\end{recap}
